We compare four communication channels in a hidden-target, two-robot rendezvous on CPU: no communication, a learned continuous message vector, a learned discrete token channel, and a fixed human-readable templated language (``target is $\langle$row word$\rangle$ $\langle$column word$\rangle$''). The templated language is only a small symbolic proxy for natural language; no language model is used. All systems are trained with the same pathwise-gradient objective through differentiable dynamics (not reinforcement learning) and evaluated over 5 seeds. Every channel beats no communication (success 0.071). The continuous vector has the highest final success (1.000), ahead of the templated language (0.931) and the discrete channel (0.678). Both the continuous and language systems are at plateau by the first evaluation (6400 episodes), so the area under the learning curve mostly reflects final success; with evaluation every 640 episodes the language reaches 50\% validation success in 1024 $\pm$ 351 training episodes versus 896 $\pm$ 351 for the continuous vector, i.e.\ no resolved speed difference, but far faster than the discrete channel (9984 $\pm$ 3282). The discrete baseline is untuned and uses only about 11 of its 16 messages; a temperature sweep over 5 seeds raises it at best to 0.721, so its position is provisional. In cross-play, listeners paired with speakers from other seeds fall to chance level for both learned channels (0.066 each, versus 0.071 for no communication), whereas the templated language keeps 0.931, which is expected because its speaker is fixed by construction rather than learned. Capacity and noise sweeps show the ranking depends strongly on grid resolution and channel noise. These results come from one toy task, a reimplemented and simplified set of baselines and 5 seeds, and should not be read as evidence about natural language.
